\section{结论与政策建议}

\subsection{主要结论}

本文利用2015—2024年中国A股上市公司面板数据，采用双向固定效应模型和手动Sobel中介效应检验，系统考察了企业数字化转型对供应链韧性的因果影响及其传导机制。主要发现如下。

第一，数字化转型对供应链韧性具有统计显著且经济意义可观的正向影响。在控制企业和年份双向固定效应后，数字化转型指数每提升一个单位，供应链韧性提高约0.0015（p=0.029），相当于DT每增加一个标准差（10.7单位），韧性提升约2\%的样本均值。这一效应在滞后解释变量和替代被解释变量函数形式两项关键稳健性检验中保持一致，排除了逆向因果关系和函数形式偏差的竞争性解释。

第二，数字化转型的供应链韧性提升效应具有鲜明的"危机放大"特征。疫情前子样本中DT系数未达显著性水平，疫情后子样本同样如此，但DT与疫情后虚拟变量的交互项系数为0.0016且在1\%水平极其显著——系统性冲击使得数字化的边际韧性保护效应放大了一倍以上。这一规律对于理解数字化基础设施的"保险"属性以及政府在危机预防阶段的数字化激励政策设计具有直接启示。

第三，效应在高财务杠杆企业和中小企业中更为集中。高杠杆企业面临更强的外部融资约束，数字化在这些企业中扮演了替代性风险管理工具的角色；中小企业数字化起步较晚、边际提升空间较大，从"数字化赋能"中获益相对于已实现高度数字化的大型企业更为明显。疫情交互项与杠杆异质性的叠加，进一步暗示了一种"双重脆弱性补偿"机制——数字化对财务脆弱和外部冲击敏感的企业群体的保护效应更大。

第四，本文最具反思性的发现在于机制检验：三条理论上合理且得到既有文献初步支持的中介路径——流程优化度、企业创新能力和信息共享水平——在严谨的Sobel检验下均未成立。这一"干净的零发现"具有三方面的学术含义。其一，它提示既有文献中基于单一中介路径的机制叙事可能需要更为审慎的检验。其二，数字化转型对供应链韧性的直接效应主导模式（全中介模型中DT系数不降反升）暗示，数字化的韧性功能可能更多通过实时监测、快速切换和精准调度等"不需要中介"的直接操作渠道实现，而非通过预设的组织能力改善间接传导。其三，信息共享水平的反向$ a $路径（DT提升却伴随信息共享下降）提示了"数据孤岛"假说的存在可能——高度数字化企业倾向于封闭内部信息体系以保护数据资产——这一问题本身值得后续研究的独立探讨。

\subsection{政策建议}

基于上述发现，本文提出以下政策建议。

第一，在数字经济政策的靶向性方面。本研究的异质性分析表明，数字化转型的供应链韧性提升效应并非在所有企业中均匀分布，而是集中在高杠杆企业和中小企业。这意味着"数字赋能实体经济"的政策口号需要转化为差异化的资源分配策略：政府对企业数字化的补贴和激励（如"上云用数赋智"行动、智能制造试点示范等）应优先向高财务脆弱性企业和中小企业倾斜，因为它们的"受益弹性"更高，单位政策投入的边际效应更大。相反，对已经实现高度数字化的龙头企业，政策的着力点应从"鼓励建设"转向"鼓励开放"——促进其将数字化基础设施以平台化方式向供应链上下游中小企业溢出。

第二，在供应链安全制度设计的时序方面。疫情交互项的显著正向发现（DT$\times$后疫情系数0.0016，p$<$0.01）传递了一个清晰的政策信号：数字化基础设施的供应链保护价值在经济平稳期往往被低估，而在系统性危机中才充分显现。这与自然灾害保险市场的"低概率事件低估"行为偏误在逻辑上同构。政策启示在于，对企业数字化的激励应当具有前瞻性和逆周期特征——在经济平稳期通过税收优惠、专项贷款和技术补贴等方式鼓励企业完成数字化基建布局，而不是等到危机暴露供应链脆弱性之后再被动补救。当前"数字经济创新发展试验区"和"供应链创新与应用试点"等政策窗口为这一"预防性数字化"思路提供了载体。

第三，在企业数字化投资策略方面。中介效应的零发现对企业管理者具有直接的决策参考价值：花钱搞数字化不会自动带来供应链韧性的提升，尤其是不会自动通过"优化流程""促进创新"或"改善信息共享"来间接增强韧性。数字化投入必须直接瞄准供应链管理的关键痛点——例如部署实时预警系统以增强冲击感知速度、建立多平台订单管理系统以实现供应渠道的快速切换、投资可视化跟踪技术以减少信息延迟——而非寄望于数字化投入发生"润物细无声"式的组织能力溢出。简言之，数字化的韧性功能是可预期的，但需要靶向投资，而非一般性的信息化升级。

\subsection{研究局限与未来方向}

本文存在四个方面的局限，也构成了后续研究的自然起点。第一，核心解释变量基于年报文本分析，其度量精度受企业年报披露策略的影响。数字化词频可能捕捉了企业"说数字化"的热情而非"做数字化"的实效。后续研究可尝试结合企业IT投资金额、数字化专利分类号和多来源数据交叉验证，构建更为综合的度量体系。第二，三条中介路径的"零发现"既可能指向理论猜想本身需要修正（传导机制可能在操作层面而非组织能力层面），也可能源于中介变量的操作化定义未能充分捕捉理论概念的核心内涵。后续研究可探索更为贴近供应链操作层面的中介变量，如供应商切换周期、订单响应时间和库存周转弹性等，以寻找数字化韧性效应的真实传导渠道。第三，因果识别主要依赖双向固定效应和滞后变量策略，尽管标准误在企业层面聚类且在滞后变量检验中排除了逆向因果，但离"完美的因果推断"仍有距离。随着更多政策冲击数据的积累（如各地数字经济试点政策的错层推进），未来的研究有望利用交错型双重差分（staggered DID）等前沿方法提供更为干净的因果识别。第四，本研究基于中国A股上市公司数据，结论对非上市中小企业和发展中国家以外的制度环境的可推广性有待进一步检验。跨制度比较——例如中国与东南亚、拉美等新兴市场国家之间的数字化—韧性关系比较——将为这一研究议程增添外部有效性的关键证据。